\documentclass[11pt]{article}
\usepackage{amssymb, amsmath, libertine, hyperref, xcolor, booktabs, enumitem, microtype}
\usepackage[T1]{fontenc}
\usepackage[margin=1in]{geometry}
\usepackage{tikz}
\newcommand{\gridxy}[4]{% xmin xmax ymin ymax -- gridlines exactly on the ticks
  \foreach \gX in {#1,...,#2} \draw[gray!25] (\gX,#3) -- (\gX,#4);
  \foreach \gY in {#3,...,#4} \draw[gray!25] (#1,\gY) -- (#2,\gY);}
\definecolor{dark-red}{rgb}{0.4,0.15,0.15}
\hypersetup{colorlinks, linkcolor=dark-red, urlcolor=dark-red}
\setlist[enumerate,1]{label=(\alph*), leftmargin=*, itemsep=5pt, topsep=5pt}
\pagestyle{empty}
\newcommand{\prob}[2]{\medskip\noindent\textbf{\large #1}\hfill\textit{#2}\par\smallskip}
\begin{document}

\begin{center}
{\Large\bfseries Math 111: Calculus \quad Practice Exam 2}\\[4pt]
{For Exam 2, in class on Friday, October 16}
\end{center}

\medskip
\noindent\textbf{What Exam 2 is like.} Four problems in a fifty-minute period, closed-book
and device-free. It is \textbf{cumulative}: the new material is Chapter~2 (the derivative
rules, including the chain rule and inverse functions), \S3.1 (related rates), and \S3.3
(extreme values), but ideas from Unit~1 --- reading $f'$ from a graph or a table,
over- and under-estimates, the tangent line --- come back inside the new problems. Each
problem is marked holistically on the 0--5 scale: a right answer with no reasoning is not a
5, and a wrong answer containing a real idea is not a 0. The derivative rules are now
fair game and expected; say which rule you are using and why it applies.

\smallskip
\noindent\textbf{How to use this.} Sit down with a blank sheet and fifty minutes on a
clock, book closed and phone in another room. Bring whatever you could not finish to
review on Wednesday the 14th. Problem~4 uses \S3.3, which we cover on Monday the 12th.

\bigskip\hrule

\prob{1. Structure first}{about 10 minutes}
For each function, first say in a few words how it is built (a sum? a product? a
quotient? a composition, and of what?), then differentiate it. You need not simplify.
\begin{enumerate}
\item $f(x) = x^2\cos x$
\item $g(x) = \dfrac{e^x}{x+1}$
\item $h(x) = \sqrt{1+\sin(3x)}$
\item $k(x) = \ln(1+x^2)$
\item Find the tangent line to the graph of $f$ at $x=\pi$, and use it to estimate
$f(3)$. (You may leave $\pi$ in your answer.)
\end{enumerate}

\prob{2. A balloon, a barometer, and two tables}{about 13 minutes}
A weather balloon rises. Its altitude $a(t)$, in meters, is recorded every five minutes:
\smallskip
\begin{center}
\begin{tabular}{@{}lrrrrr@{}}
\toprule
$t$ (minutes) & 0 & 5 & 10 & 15 & 20\\
$a(t)$ (m) & 0 & 400 & 750 & 1050 & 1300\\
\bottomrule
\end{tabular}
\end{center}
\smallskip
Air pressure $P(a)$, in kilopascals, depends only on altitude $a$:
\smallskip
\begin{center}
\begin{tabular}{@{}lrrrrr@{}}
\toprule
$a$ (m) & 0 & 500 & 1000 & 1500 & 2000\\
$P(a)$ (kPa) & 101.3 & 95.5 & 90.0 & 84.8 & 79.9\\
\bottomrule
\end{tabular}
\end{center}
\smallskip
The pressure the balloon experiences at time $t$ is $H(t)=P\big(a(t)\big)$. There are no
formulas; the tables are all you have.
\begin{enumerate}
\item Estimate $a'(10)$, with units. Then give one estimate of $a'(10)$ that is certainly
too large and one that is certainly too small, and say what you are assuming about $a'$
in order to be sure.
\item Estimate $P'(750)$, with units. Say in one sentence what the sign of this number
means physically.
\item Estimate $H'(10)$, with units. Explain why the two numbers you combine should be
combined the way you combined them --- in terms of what happens to the balloon over the
next minute, not by naming a rule.
\item A classmate computes $H'(10)$ as $P'(10)\cdot a'(10)$. Say precisely what is wrong,
and why the number $P'(10)$ is not even the kind of quantity that belongs here.
\end{enumerate}

\prob{3. A melting snowball}{about 12 minutes}
A snowball is melting. Assume it stays a perfect sphere, and that its volume is
decreasing at a steady $2$ cubic centimeters per minute. Recall $V=\frac43\pi r^3$ and
surface area $S=4\pi r^2$.
\begin{enumerate}
\item Draw and label a picture, and write down, with units, every rate you know and the
rate you want.
\item How fast is the radius shrinking at the moment the radius is $3$ cm?
\item How fast is the surface area shrinking at that moment?
\item Is the radius shrinking faster or slower when the snowball is smaller? Answer from
your formula in (b), then say in a sentence why that makes physical sense.
\end{enumerate}

\prob{4. All about $f'$}{about 13 minutes}
A function $f$ is differentiable everywhere, and its derivative is
\[
f'(x) = (x-2)(x+1)^2 .
\]
You are not told $f$ itself.
\begin{enumerate}
\item Find all critical numbers of $f$.
\item Make a sign chart for $f'$, and use the first derivative test to classify each
critical number as a local maximum, a local minimum, or neither. Justify each from the
chart.
\item Compute $f''$ and apply the second derivative test at each critical number. Where
the test is silent, say so, and say which test settles the matter.
\item A classmate writes: ``$f'(-1)=0$, and $f'$ is negative on both sides of $-1$, so the
graph of $f$ turns around at $x=-1$ and $f$ has a local maximum there.'' Find the first
wrong step, say why it is wrong, and describe what the graph of $f$ actually does near
$x=-1$.
\item On axes of your own, sketch a possible graph of $f$ on $[-3,4]$, consistent with
everything above. Mark the critical points.
\end{enumerate}

\end{document}
